\section{Operation Modes}\label{operation-modes}

The synth provides different operation modes, each optimised for a
different task. The active mode changes the function of the user
interface controls.

\subsection{Voice Edit Mode}\label{voice-edit-mode}

The voice edit mode is used to modify sound parameters and to do basic
pattern editing.

\begin{itemize}
\tightlist
\item
  \textbf{Display (1):} parameters of the selected synthesis page.
\item
  \textbf{Knobs (2):} edit the visible parameters.
\item
  \textbf{Sequencer buttons (5):} set and remove the 16 main steps of
  the current track.
\item
  \textbf{Select buttons (7):} select the active synthesis page
  (oscillator, envelope, mixer, etc.)
\item
  \textbf{Voice buttons (9):} select the voice and sequencer track to
  edit; with shift, mute single voices.
\end{itemize}

\subsubsection{Pattern Editing}\label{pattern-editing}

The 16 sequencer buttons set and remove steps in the selected pattern.
Each button is a \(\frac{1}{16}\)th note from a 4/4 bar. A lit LED above
the button indicates the step is set.

\subsubsection{Step Parameter Editing}\label{step-parameter-editing}

\begin{enumerate}
\def\labelenumi{\arabic{enumi}.}
\tightlist
\item
  Press and hold the shift button.
\item
  The display shows the step parameter page of the selected step.
\item
  Use the knobs and encoder to change parameters.
\item
  A flashing LED indicates the selected step.
\item
  Select a different step using the 16 sequencer buttons.
\end{enumerate}

\subsubsection{Sub Step Editing}\label{sub-step-editing}

Sub steps can be set and removed in voice edit mode:

\begin{enumerate}
\def\labelenumi{\arabic{enumi}.}
\tightlist
\item
  Press and hold the shift button.
\item
  Use the 16 sequencer buttons to select the main step whose sub steps
  you want to edit.
\item
  A flashing LED indicates the selected step.
\item
  The sub step pattern is shown on the 8 select button LEDs.
\item
  Set and remove sub steps using the 8 select buttons.
\end{enumerate}

\attention{The step parameter menu always shows the parameters of the last-tweaked step.
Setting or removing a sub step causes the display to show that sub step's parameters. For more
ergonomic step parameter editing, use Step (Edit) mode.}

\subsubsection{Sound Editing}\label{sound-editing}

\begin{itemize}
\tightlist
\item
  The 7 voice buttons select the active voice (lit LED indicates
  selection).
\item
  Voice parameters are divided into pages selectable with the 8 select
  buttons.
\item
  Each voice type has its own parameter set, but the overall page
  structure is the same.
\item
  Parameters shown in the display are edited with the 4 knobs.
\end{itemize}

\textbf{Example --- changing the filter frequency of voice 3:}

\begin{enumerate}
\def\labelenumi{\arabic{enumi}.}
\tightlist
\item
  Press Voice Button 3 to select the voice.
\item
  Press Select Button 6 to show the filter page.
\item
  Use the first knob to adjust the \param{frq} parameter.
\end{enumerate}

\subsection{Performance Mode}\label{performance-mode}

This mode is designed for playing songs live. You can chain and change
patterns on the fly, trigger manual rolls for each voice, and access the
morph and global sample rate effects. Voice editing is not available in
this mode.

\begin{figure}[H]
  \centering
  \includegraphics[width=\linewidth]{performance_mode_panel.jpeg}
  \caption{LXR front panel in performance mode.}
\end{figure}

\begin{itemize}
\tightlist
\item
  \textbf{Display (1):} performance parameters.
\item
  \textbf{Knobs (2):} change the performance parameters.
\item
  \textbf{Sequencer buttons 1--7 (5):} trigger a manual roll for each
  corresponding voice.
\item
  \textbf{Select buttons (7):} change the playing pattern.
\item
  \textbf{Voice buttons (9):} mute/unmute voices without holding shift.
\end{itemize}

\subsubsection{Basic Performance Menu}\label{basic-performance-menu}

\begin{paramtable}
rol & Manual roll rate  & Trigger rate for the manual roll function of sequencer buttons 1–7. \\
mrp & Morph amount      & Ratio between the original sound and the morph target. \\
sr  & Global sample rate & A global sample rate reduction effect. \\
shu & Shuffle amount    & Sets the amount of global shuffle. \\
\end{paramtable}

\subsubsection{Manual Roll}\label{manual-roll}

The first 7 sequencer buttons trigger a manual roll for each
corresponding voice. The roll rate is set with the \param{rol}
parameter. Rolls are recorded to the pattern when recording mode is
active.

\subsubsection{Morph}\label{morph}

The morph feature transitions from one preset sound to another. Load any
preset from the SD card as a morph target. The \param{mrp} parameter
controls the ratio between the original and target sound. As morph
increases, the current sound is gradually transformed into the morph
target.

For loading a morph target, see the Load and Save Mode section.

\subsubsection{Shuffle}\label{shuffle}

The shuffle function shifts the position of every other
\(\frac{1}{16}\)th note.

Original timing:
\texttt{\textbar{}\ X\ x\ x\ x\ X\ x\ x\ x\ X\ x\ x\ x\ X\ x\ x\ x\ \textbar{}}

With shuffle applied:
\texttt{\textbar{}\ X\ xx\ xX\ xx\ xX\ xx\ xX\ xx\ x\ \textbar{}}

\subsubsection{Change Pattern}\label{change-pattern}

The 8 select buttons change the playing pattern, selecting from the 8
patterns of the loaded pattern set. Pattern changes occur only at the
end of a bar (track 1 is the reference track when different track
lengths are used).

\subsubsection{Chain Patterns}\label{chain-patterns}

To build longer pattern chains, go to the pattern options screen:

\begin{enumerate}
\def\labelenumi{\arabic{enumi}.}
\tightlist
\item
  Press and hold shift.
\item
  A flashing LED indicates the selected pattern.
\item
  The display shows pattern options.
\end{enumerate}

\oled{Rpt \quad nxt}{1 \hspace{2.5em} p2}

For each pattern you can specify the \textbf{next pattern} to play and
the \textbf{repeat/change measure}.

\textbf{Change Measure / Repeat (rpt):} The change measure is relative
to the absolute song position in bars. This ensures pattern chains stay
in sync regardless of when you manually switch patterns.

\textbf{Next Pattern (nxt):} Values p1--p8 select a specific pattern.
Values r2--r8 select a random pattern between pattern 1 and the selected
value (e.g.~r4 randomly selects from patterns 1--4).

\begin{figure}[H]
  \centering
  \includegraphics[width=\linewidth]{pattern_chain_example1.png}
  \caption{Pattern chain example: pattern 1 plays 3 times, then pattern 2 plays once.}
\end{figure}

\begin{figure}[H]
  \centering
  \includegraphics[width=\linewidth]{pattern_chain_example2.jpeg}
  \caption{Pattern chain with manual switch: the break always lands on bar 4.}
\end{figure}

\tip{When editing multiple chained patterns, set the pattern follow (\param{flw}) parameter in the
settings menu to `Off`. This lets you edit a pattern without it jumping to follow playback.}

\subsubsection{Follow Mode and Viewed
Pattern}\label{follow-mode-and-viewed-pattern}

The follow mode setting is in the global settings menu (on by default).

When follow mode is active, the viewed pattern always matches the
playing pattern --- all LEDs and display values update automatically on
every pattern change.

When follow mode is off, the user must manually select the viewed
pattern. This allows editing a different pattern from the one currently
playing.

\textbf{Selecting the viewed pattern (performance mode only):}

\begin{enumerate}
\def\labelenumi{\arabic{enumi}.}
\tightlist
\item
  The currently viewed pattern blinks; the playing pattern shows a solid
  LED.
\item
  Hold shift.
\item
  Use the 8 select buttons to pick the pattern to view/edit.
\item
  Release shift.
\end{enumerate}

\subsection{Step (Edit) Mode}\label{step-edit-mode}

Step edit mode is for editing step parameters more efficiently. Think of
it as an inverted voice edit mode with respect to the sequencer.

The display always shows the step parameter page. The sequencer and
select buttons are only used to select steps --- you can edit step
parameters without holding shift.

The 16 sequencer buttons select a main step and show its 8 sub steps on
the select button LEDs. The selected step is indicated by a blinking
LED.

\subsubsection{Setting and Removing
Steps}\label{setting-and-removing-steps}

To set or remove steps from the pattern, hold the shift button while
pressing the sequencer buttons.

\subsection{Load and Save Mode}\label{load-and-save-mode}

The 4th mode button opens the load/save page. Pressing it again toggles
between load and save.

A push on the Load/Save button brings up the load page:

\oled{Load: Sound}{[ \quad 1] Name}

Pressing the button again shows the save page:

\oled{Save: Sound}{[ \quad 1] Name \hfill ok}

\subsubsection{Display Description}\label{display-description}

\begin{itemize}
\tightlist
\item
  \textbf{Row 1:} active mode (Load or Save) and the selected data type
  (Sound, Pattern, Morph Sound, Settings).
\item
  \textbf{Row 2:} active preset number and the name of the selected
  preset.
\item
  Pages with an \texttt{ok} button only execute when you navigate to
  \texttt{ok} and push the encoder.
\item
  Pages without an \texttt{ok} button execute as soon as a change
  occurs.
\end{itemize}

\subsubsection{Menu Navigation}\label{menu-navigation}

Use the encoder to navigate and change entries. Square brackets
\texttt{{[}\ {]}} show the field that will be changed by the encoder
(edit mode). Pushing the encoder switches to selection mode (shows
\texttt{\textgreater{}}).

\tip{For the data type field (Sound, Pattern, Morph Sound, Settings), use the first knob to select
directly without clicking through encoder modes.}

\subsubsection{Sound Presets}\label{sound-presets}

Sound presets contain all synthesis parameters for the 6 voices.

\textbf{Loading:} sounds load automatically when the preset number
changes.

\textbf{Saving:}

\begin{enumerate}
\def\labelenumi{\arabic{enumi}.}
\tightlist
\item
  Select the preset number.
\item
  Navigate to the name letters.
\item
  Push the encoder and turn it to change each letter.
\item
  Navigate to \texttt{ok} and push the encoder to save.
\end{enumerate}

\textbf{Quick naming:} while the cursor is in the name area, knob 1
moves the cursor, knob 2 toggles case, and knob 3 scrolls through ASCII
characters.

\subsubsection{Morph Presets}\label{morph-presets}

A morph preset is a normal preset loaded as a morph target.

\textbf{Loading:}

\begin{enumerate}
\def\labelenumi{\arabic{enumi}.}
\tightlist
\item
  Select \texttt{MorphSnd} as the data type.
\item
  Change the preset number --- the morph sound loads automatically.
\item
  Note: no audible change occurs if the \param{mrp} (morph amount)
  parameter is set to zero.
\end{enumerate}

\textbf{Saving morph results:} select \texttt{MorphSnd} on the save
screen to save the current blended sound to a new preset slot.

\subsubsection{Pattern Sets}\label{pattern-sets}

Pattern sets are collections of 8 patterns stored in one file.

\textbf{Loading:} navigate to \texttt{OK} on the display and push the
encoder to start loading (required because loading takes 4--5 seconds).
The pattern loads in the background --- \textbf{the loaded pattern does
not play until you press one of the pattern change buttons in
performance mode.}

\textbf{Saving:} same procedure as saving a sound preset.

\subsubsection{Settings}\label{settings}

The settings file contains all parameters from the settings menu. It is
automatically loaded at boot. To save current settings, go to save mode,
select \texttt{Globals}, navigate to \texttt{Ok}, and push the encoder.

\subsection{Pattern Generator Mode}\label{pattern-generator-mode}

The pattern generator creates polyrhythms automatically using the
Euclidean algorithm.

\subsubsection{Accessing the Pattern
Generator}\label{accessing-the-pattern-generator}

\begin{enumerate}
\def\labelenumi{\arabic{enumi}.}
\tightlist
\item
  Press and hold shift.
\item
  Press Mode Button 2 (Perf/Pattern).
\item
  A flashing mode 2 LED indicates the pattern generator is active.
\end{enumerate}

\subsubsection{Pattern Generator Menu}\label{pattern-generator-menu}

\begin{paramtable}
len & Track Length   & Length of the track in $\frac{1}{16}$ notes. Each track can have a
                       different length. \\
stp & Number of Steps & The number of active steps to distribute in the pattern. \\
\end{paramtable}

\subsubsection{Generating Patterns}\label{generating-patterns}

\begin{itemize}
\tightlist
\item
  The generator modifies a single track at a time. Select the track with
  the 7 voice buttons.
\item
  Changing either \param{len} or \param{stp} immediately generates a new
  pattern on the active track.
\end{itemize}

\attention{Existing pattern data is overwritten and there is no undo.}

\subsection{Global Settings Menu}\label{global-settings-menu}

\subsubsection{Accessing the Settings
Menu}\label{accessing-the-settings-menu}

\begin{enumerate}
\def\labelenumi{\arabic{enumi}.}
\tightlist
\item
  Press and hold shift.
\item
  Press Mode Button 4 (Load/Save).
\item
  A flashing mode 4 LED indicates the settings menu is active.
\end{enumerate}

\subsubsection{Menu Options}\label{menu-options}

\begin{paramtable}
bpm & Tempo in BPM    & Internal clock tempo. Set to 0 to sync to incoming MIDI clock. \\
fet & Parameter fetch & When active, you must `fetch' the displayed value with the knob
                        before it changes — prevents value jumps on page switches. \\
flw & Pattern follow  & When active, the viewed pattern always matches the playing pattern. \\
qnt & Sequencer Quantisation & Sets the quantisation grid for recorded data. \\
ss  & Screensaver     & Turns the screensaver on or off. \\
\end{paramtable}

\subsubsection{Follow Mode}\label{follow-mode}

When follow mode is active, all viewed pattern data (step LEDs, display
values) updates automatically on every pattern change. When off, the
user selects the viewed pattern manually --- allowing editing of a
pattern other than the one playing. The chaselight is only visible when
the playing pattern is viewed.

\subsubsection{Screensaver}\label{screensaver}

The display is an OLED. To extend its lifetime and prevent burn-in, a
screensaver activates after a few minutes of inactivity. Any control
touch dismisses it. The screensaver can be disabled in the settings
menu.

\subsubsection{Loading and Saving Global
Settings}\label{loading-and-saving-global-settings}

Settings are loaded automatically on power-up. To save current settings,
go to the save menu, select \texttt{Globals}, navigate to \texttt{Ok},
and push the encoder.
